\section{Conclusiones}

Se concluye exitoso la optimizacion del algoritmo original presentado con los pasos desarrollados en este informe, transformadose de un algoritmo de O(n*n) a O(n*log(n)). Se
deja expresado tambien que no es la mejor optimizacion que puede obtenerese para este algoritmo debido a que no se realizó una investigacion exhaustiva en metodologias 
 o tecnicas que permitar micro optimizar estas dos acciones del algoritmo que son *hallar suma de divisiores* y *comparar numeros*. 

Se deja tambien expresado que la idea de utilizar una criba para la optimizacion de este elgoritmo fué obtenida al discutir con previos estudiantes de la materia Teoria de 
Algoritmos de la Universidad de Buenos Aires, entrevistados durante el desarrollo de este Trabajo Practico.

Otra mencion a realizar es la utilizacion de herramientas de inteligencia artificial para el desarrollo de este Trabajo. Principalmente para la creacion de los entornos de 
pruebas de algoritmos con el fin de obtener los graficos ya presentados y la redaccion de estas paginas en el lenguaje LaTex con el compilador LatexMk 4.88 debido a mi escasa
experiencia practica con este. 

